%------------------------------------------------------------------------------%
\documentclass[fleqn,utf8,aspectratio=169,12pt]{beamer}

\usepackage{integracao_e_series}

\title{Teorema Fundamental do Cálculo -- Parte 2}

%------------------------------------------------------------------------------%
\begin{document}

\addtitle

%------------------------------------------------------------------------------%
\begin{frame}{Teorema Fundamental do Cálculo -- Parte 2}
  Se $f$ for {\color{blue} contínua} em $[a, b]$
  \pause
  então, para {\color{blue}qualquer primitiva} $F$ de $f$,
  \pause
  \[
    \int_a^b f(x) \,dx      \pause
    \;=\; F(b) - F(a)       \pause
    \;=\; F(x) \evalat{a}{b}
  \]
\end{frame}

%------------------------------------------------------------------------------%
\begin{proofframe}{Demonstração}
  Queremos mostrar que
  \[
    \int_a^b f(x) \,dx = F(b) - F(a)
  \]
\end{proofframe}

%------------------------------------------------------------------------------%
\begin{proofframe}{Demonstração}
  Pelo Teorema Fundamental do Cálculo -- Parte 1
  \[
    G(x) = \int_a^x f(t)\,dt
  \]
  é uma primitiva de $f$

  \pause
  \bigskip
  Qualquer outra primitiva tem a forma
  \[
    F(x) = G(x) + C
  \]
\end{proofframe}

%------------------------------------------------------------------------------%
\begin{proofframe}{Demonstração}
  \begin{align*}
    \onslide<+->{F(b) - F(a)}
    \onslide<+->{& = \big(G(b) + C\big) - \big(G(a) + C\big) \\[5mm]}
    \onslide<+->{& = G(b) - G(a)                             \\[3mm]}
    \onslide<+->{& = \int_a^b f(x)\,dx - \int_a^a f(x)\,dx   \\[3mm]}
    \onslide<+->{& = \int_a^b f(x)\,dx}
  \end{align*}
\end{proofframe}

%------------------------------------------------------------------------------%
\section{Exemplos}

%------------------------------------------------------------------------------%
\begin{question}
  Calcular a ``área'' sob o gráfico de \(f(x)=x^2\) entre 0 e $b$
\end{question}

%------------------------------------------------------------------------------%
\begin{resolution}{Solução}
  A ``área'' sob o gráfico é
  \[
    \pause
    A
    \;=\; \int_a^b f(x) \,dx
    \pause
    \;=\; \int_0^b x^2  \,dx
  \]
\end{resolution}

%------------------------------------------------------------------------------%
\begin{resolution}{Solução}
  Calculamos a primitiva de $f(x)=x^2$
  \[
    F(x)
    \;=\; \int x^2 \,dx
    \pause
    \;=\; \frac{x^{2+1}}{2+1} + C
    \pause
    \;=\; \frac{x^3}{3} + C
  \]
\end{resolution}

%------------------------------------------------------------------------------%
\begin{resolution}{Solução}
  Pelo Teorema Fundamental do Cálculo,
  \[
    A
    \;=\; F(x) \evalat{0}{b}
    \pause
    \;=\; F(b) - F(0)
    \pause
    \;=\; \left[ \frac{b^3}{3} + C \right]
      -   \left[ \frac{0^3}{3} + C \right]
    \pause
    \;=\; \frac{b^3}{3}
  \]
\end{resolution}

%------------------------------------------------------------------------------%
%------------------------------------------------------------------------------%
\begin{question}
  Calcule \(\int_3^6 \frac{dx}{x}\)
\end{question}

%------------------------------------------------------------------------------%
\begin{resolution}{Solução}
  \(f(x)=\frac{1}{x}\) é contínua no intervalo $[3, 6]$

  \pause
  \bigskip
  Uma primitiva de $f$ é $F(x)=\ln\abs{x}$

  \pause
  \bigskip
  Pelo Teorema Fundamental do Cálculo
  \[
    \int_3^6 \frac{dx}{x}
    \pause
    = \ln\abs{x} \;\evalat{1}{3}
    \pause
    = \ln(6) - \ln(3)
    \pause
    = \ln\left(\frac{6}{3}\right)
    \pause
    = \ln(2)
    \pause
    \approx 0.693
  \]
\end{resolution}

%------------------------------------------------------------------------------%
%------------------------------------------------------------------------------%
\begin{question}
  Calcule
  \(
    \int_{-\pi}^\pi \sin(x) \; dx
  \)
\end{question}

%------------------------------------------------------------------------------%
\begin{resolution}{Solução}
  Queremos uma \emph{integral definida} da função
  \[
    f(x) = \sin(x)
  \]

  \pause
  Primeiro precisamos encontrar uma primitiva de $f$

  \pause
  \bigskip
  Tendo a primitiva, usamos o Teorema Fundamental do Cálculo
\end{resolution}

%------------------------------------------------------------------------------%
\begin{resolution}{Primitiva}
  \begin{align*}
    \onslide<+->{F(x) & = \int f(x) dx    \\[2mm]}
    \onslide<+->{     & = \int \sin(x) dx \\[2mm]}
    \onslide<+->{     & = -\cos(x) + c}
  \end{align*}
\end{resolution}

%------------------------------------------------------------------------------%
\begin{resolution}{Solução}
  \begin{align*}
    \onslide<+->{I}
    \onslide<+->{& = \int_{-\pi}^\pi \sin(x) dx }
    \onslide<+->{\;=\; F(x)\evalat{-\pi}{\pi}}
    \onslide<+->{\;=\; F(-\pi) - F(\pi)}\\
    \onslide<+->{& = \left(-\cos(x)\right) \evalat{-\pi}{\pi}\\}
    \onslide<+->{& = -\cos(\pi)-\big(-\cos(-\pi)\big)\\}
    \onslide<+->{& = -\cos(\pi)+\cos(-\pi)\\}
    \onslide<+->{& = -(-1)+(-1)\\}
    \onslide<+->{& = 0}
  \end{align*}
\end{resolution}

%------------------------------------------------------------------------------%
%------------------------------------------------------------------------------%
\begin{question}
  Calcule
  \(
    \int_1^x x\cos(\pi t) \, dt
  \)

  \pause
  \bigskip
  \bigskip
  Note que a integral é em $t$ portanto $x$ é uma constante
\end{question}

%------------------------------------------------------------------------------%
\begin{resolution}{Solução}
  Queremos calcular
  \[
    A
    \;=\; \int_1^x x\cos(\pi t) \, dt
    \pause
    \;=\; x \int_1^x \cos(\pi t) \, dt
  \]

  \pause
  Calculando a primitiva
  \[
    \pause
    F(t)
    \pause
    \;=\; \int \cos(\pi t) \, dt
    \pause
    \;=\; \frac{\sin(\pi t)}{\pi} + c
  \]
 \end{resolution}

%------------------------------------------------------------------------------%
\begin{resolution}{Solução}
\begin{align*}
  \onslide<+->{A}
  \onslide<+->{& = x \, \big(F(t)\big) \evalat{1}{x}                                   \\[1mm]}
  \onslide<+->{& = x \, \left(\frac{\sin(\pi t)}{\pi}\right) \evalat{1}{x}             \\[1mm]}
  \onslide<+->{& = x \, \left(\frac{\sin(\pi x)}{\pi} - \frac{\sin(\pi 1)}{\pi}\right) \\[1mm]}
  \onslide<+->{& = x \, \left(\frac{\sin(\pi x)}{\pi} - 0\right)                       \\[1mm]}
  \onslide<+->{& = \frac{x}{\pi}\sin(\pi x)}
\end{align*}
\end{resolution}


%------------------------------------------------------------------------------%
%------------------------------------------------------------------------------%
\begin{question}
  Calcule
  \(
    \int_0^1 20x^5 - e^x + 3x^6 - \cos(x) + \frac{3}{1+x^2} - 3 \; dx
  \)
\end{question}

%------------------------------------------------------------------------------%
\begin{resolution}{Solução}
  Queremos uma \emph{integral definida} da função
  \[
    f(x) = 20x^5 - e^x + 3x^6 - \cos(x) + \frac{3}{1+x^2} - 3
  \]

  \pause
  Primeiro precisamos encontrar uma primitiva de $f$

  \pause
  \bigskip
  Tendo a primitiva, usamos o Teorema Fundamental do Cálculo
\end{resolution}

%------------------------------------------------------------------------------%
\begin{resolution}{Solução}
  \begin{align*}
    \onslide<+->{F(x) & = \int f(x) dx \\[2mm]}
    \onslide<+->{& = \int 20x^5 - e^x + 3x^6 - \cos(x) + \frac{3}{1+x^2} - 3\; dx \\[2mm]}
    \onslide<+->{& = \int 20 x^5 dx
      - \int e^x dx
      + \int 3 x^6 dx
      - \int \cos(x) dx
      + \int \frac{3}{1+x^2} dx
      - \int 3 dx \\[2mm]}
    &
      \onslide<+->{= 20\frac{x^6}{6}}
      \onslide<+->{- e^x}
      \onslide<+->{+ 3\frac{x^7}{7}}
      \onslide<+->{- \sin(x)}
      \onslide<+->{+ 3\arctan(x)}
      \onslide<+->{- 3x }
      \onslide<+->{+ C} \\[2mm]
    &
      \onslide<+->{
      = \frac{10}{3}x^6
      - e^x
      + \frac{3}{7}x^7
      - \sin(x)
      + 3\arctan(x)
      - 3x
      + C}
  \end{align*}
\end{resolution}

%------------------------------------------------------------------------------%
\begin{resolution}{Solução}
  \begin{align*}
    \onslide<+->{I}
    \onslide<+->{& = \int_0^1 f(x) dx }
    \onslide<+->{  \;=\; F(x)\evalat{0}{1}}
    \onslide<+->{  \;=\; F(1) - F(0)}
    \\
    \onslide<+->{& = \left(
        \frac{10}{3}x^6
      - e^x
      + \frac{3}{7}x^7
      - \sin(x)
      + 3\arctan(x)
      - 3x
      \right) \evalat{0}{1}
    \\}
    \onslide<+->{& = \left(
        \frac{10}{3}\times 1
      - e^1
      + \frac{3}{7}\times 1
      - \sin(1)
      + 3 \arctan(1)
      - 3\times 1
        \right)
    \\}
    \onslide<+->{& - \left(
          \frac{10}{3}\times 0
        - e^0
        + \frac{3}{7}\times 0
        - \sin(0)
        + 3 \arctan(0)
        - 3\times 0
        \right)
    \\}
  \end{align*}
\end{resolution}

%------------------------------------------------------------------------------%
\begin{resolution}{Solução}
  \begin{align*}
    \onslide<+->{I & = \left(
        \frac{10}{3}\times 1
      - e^1
      + \frac{3}{7}\times 1
      - \sin(1)
      + 3 \arctan(1)
      - 3\times 1
        \right)
    \\
    & - \left(
          \frac{10}{3}\times 0
        - e^0
        + \frac{3}{7}\times 0
        - \sin(0)
        + 3 \arctan(0)
        - 3\times 0
        \right)
    \\}
    \onslide<+->{& = \left(\frac{10}{3}}
    \onslide<+->{-e}
    \onslide<+->{+\frac{3}{7}}
    \onslide<+->{-\sin(1)}
    \onslide<+->{+ \frac{3\pi}{4}}
    \onslide<+->{ - 3 \right)}
    \onslide<+->{ - (-1)}
    \\
    \onslide<+->{& = \frac{10}{3}
                   - e
                   + \frac{3}{7}
                   - \sin(1)
                   + \frac{3\pi}{4}
                   - 2}
  \end{align*}
\end{resolution}

%------------------------------------------------------------------------------%

%------------------------------------------------------------------------------%
\section{Lista Mínima}

%------------------------------------------------------------------------------%
\begin{frame}{Lista Mínima}
  Estudar a Seção 2.5 da Apostila

  Fazer os exercícios: 1a-f, 2, 6a-f, 7, 9a

  \vfill
  Atenção: A prova é baseada no livro, não nas apresentações
\end{frame}

%------------------------------------------------------------------------------%
\end{document}
%------------------------------------------------------------------------------%
